\section{Discussion}\label{sec:discussion}
The results support the three claims. Asynchronous execution helps most where it should: under a persistent straggler and under runtime heterogeneity, the generation barrier is the dominant cost and removing it preserves throughput that the synchronous baseline loses. For homogeneous, near-instantaneous simulators at large worker counts the asynchronous coordination overhead does not pay off: the multi-node Lotka--Volterra points, where the fair synchronous baseline overtakes the asynchronous method, make this honest boundary explicit. On the realistic Cellular Potts workload (the cost-bearing regime the design targets) the method scales almost linearly while the synchronous baseline, held at the generation barrier even when given a matching population, scales only sub-linearly --- at a small measured cost in posterior quality on that benchmark ($\S$\ref{sec:results-validity}), on a metric that penalizes the concentration the extra throughput buys.

\section{Limitations}\label{sec:limitations}
We are explicit about what \S\ref{sec:theory} does and does not give. (i) AMIS estimators are consistent but not finite-sample unbiased, since proposals depend on past particles; a local-decision rule in the spirit of \cite{paige2014cascade} that recovers finite-sample unbiasedness for ABC is future work. (ii) The bandwidth floor $\epsilon_\infty>0$ in Assumption~\ref{ass:stab} means the theorems characterize the smooth-ABC posterior at the limit bandwidth, not the exact posterior at $\epsilon=0$; and for the central limit theorem (though not for consistency) that limit must be \emph{deterministic}, since the centred integrand it is built on would otherwise vary with the run. Setting \texttt{min\_tol} bounds the limit below but only a schedule that reaches its floor makes it deterministic. The reported runs leave \texttt{min\_tol} unset and use the data-driven schedule, so neither holds, and the bandwidth a finite budget reaches is not an asymptotic limit at all. (iii) The theorems hold up to the limiting ratio $r$ of \eqref{eq:fidelity} between the density the sample was drawn from and the function the estimator divides by; $r\equiv1$ gives the smooth-ABC posterior exactly, and \S\ref{sec:theory-r} decomposes $r$ into four factors and \emph{measures two of them} --- the prior floor and the snapshot quadrature --- at $\widehat\zeta=0.022$ on the Cellular Potts production run, a $1.1\%$ shift in total variation against a large-$m$ reference. The other two, the truncation constant and the capped-rejection/redraw conditioning, are bounded rather than measured, and because the factors multiply that is a partial accounting rather than a bound on the whole. We do not prove $r\equiv1$. Under asynchronous execution one of the four factors additionally absorbs the difference between the emitted proposal density and the conditional law under a filtration the proofs can use (Assumption~\ref{ass:filtration}), which is a property of the schedule rather than of the estimator. The stabilization clause (Assumption~\ref{ass:stab}(ii)) is not merely unproved for the top-$k$ archive rule but measurably violated in the configuration we run: the proposal's distance to its limit decays far too slowly (\S\ref{sec:theory-r}), so the central limit theorem covers the damped- or frozen-adaptation variants rather than these experiments, while consistency, which needs no rate, is unaffected. (iv) Keeping every accepted particle (unlike pyABC's discard-latecomers rule) over-represents fast-simulator parameter regions in the raw archive. The parameter-coupled-runtime study (\S\ref{sec:results-hpc}, Fig.~\ref{fig:param-bias}) quantifies this: steepening the coupling cuts throughput as the fast region comes to dominate the raw counts, yet --- across the sweep, a steeper regime, and with or without AMIS --- neither the unweighted posterior mean nor the AMIS-reweighted posterior's mean drifts with the coupling (AMIS reshapes the tails, not the mean), and a drain-after-deadline test finds the completion-time censoring negligible in effect at these levels (numbers there). We therefore do not correct for it algorithmically and do not claim the AMIS reweighting cancels it: dividing by the proposal density corrects proposal adaptation, not selection by completion time, so a formal correction would need completion probabilities or an anytime-sampling argument, and a residual finite-sample effect at stronger couplings cannot be ruled out --- both left to future work. (v) The snapshot buffer is fixed at $S=20$; adaptive sizing is a natural extension. (vi) Simulation-based calibration of the reported posterior --- the full-history AMIS estimator \eqref{eq:posterior-estimator} --- shows it well calibrated in one dimension and robust to multimodality (retaining both modes of a bimodal target in ${\approx}86\%$ of trials), and conservative rather than over-confident on the four-dimensional g-and-k model, where it mildly \emph{over}-covers. This four-dimensional behaviour is a reporting-support effect, not a weighting or budget failure (\S\ref{sec:results-validity}): the efficient top-$k$ archive used for online proposal is over-concentrated and would under-cover, the full evaluated history is over-dispersed and over-covers, and coverage rises monotonically through nominal as the reported support grows from one to the other, so the calibrated posterior is bracketed between the two. We report the conservative full-history end by default; an intermediate reported support or a larger archive would tighten the four-dimensional report to nominal. The denominator's prior floor acts in the same direction for a different reason, and the two are worth not conflating: the floor suppresses the far tails of the reported posterior (\S\ref{sec:theory-r}(a)) and so tightens it, whereas widening the reported support adds tail mass. Both are small next to the effect measured here --- the floor and the snapshot quadrature together move the reported posterior by about one percent in total variation (\S\ref{sec:theory-r}). (vi-b) On an identified target the reported posterior's effective sample size is set by the archive size and not by the simulation budget: $\mathrm{ESS}/k \in [1.9, 3.6]$ across a sweep of $k$ and $n$, and $237$--$303$ effective particles out of $1.1$--$1.3\times10^6$ evaluated on the production Gaussian run (\S\ref{sec:results-validity}). The bound is on \emph{identified} targets specifically --- the confounded Cellular Potts target retains an ESS fraction of $0.21$, because its bandwidth never tightens enough to concentrate the weights --- so it binds exactly where the posterior is worth having. Barrier-free execution therefore delivers more \emph{evaluations} per second, which is what \S\ref{sec:results-hpc} measures, but converting those into posterior accuracy is limited by a reporting rule we did not design for that purpose. Raising $k$ raises the effective sample proportionally and is the obvious lever; a bandwidth rule that does not track the archive downwards, or reporting at a deliberately looser $\epsilon$ than the tightest reached, are the others. We have not swept any of them for this purpose, and until that is done the honest scope of the systems result is throughput rather than time-to-accuracy. (vii) The synchronous baseline is matched on the ABC kernel and bandwidth schedule, but it still differs from our method in more than the barrier: it adds a probabilistic-rejection step, its generation length depends on the resulting acceptance probability, and it retains classical population weighting rather than the streaming estimator. Matching the kernel removes the acceptance \emph{rule} as a confound, and the \emph{barrierized twin} of \S\ref{sec:results-hpc} (Table~\ref{tab:twin}) closes the remaining gap on the straggler benchmark: the same propagator with a collective barrier before each breed runs $21$--$402\times$ slower than the asynchronous arm, so on that workload the throughput gain is attributable to synchronization rather than merely dominated by it. The twin does \emph{not} isolate synchronization for the posterior comparison, only for the throughput one: a collective barrier needs identical per-rank call counts, so the twin cannot be stopped on wall-clock and the two arms end at budgets that differ by two orders of magnitude (\S\ref{sec:results-hpc}). A matched-evaluation-count twin, which would separate the barrier's effect on the importance weights from the effect of the evaluations it forgoes, is the obvious next control and we have not run it. The granularity objection --- that the twin synchronizes every $W$ evaluations, a finer generation than the baseline's population of $100$ --- is measured rather than assumed away: re-running the twin with the barrier coarsened to the baseline's population changes its throughput by less than $1\%$ at every slowdown from $5\times$ up, and matters ($2.3$--$3.3\times$) only at the controls where barrier latency and not straggler waiting is the bottleneck (Table~\ref{tab:twin}). The twin has since been repeated under runtime heterogeneity (Table~\ref{tab:twin-hetero}), where it costs $1.1\times$ at the $\sigma{=}0$ control and $25.3\times$ at $\sigma{=}2$ at the baseline-matched granularity, so the attribution now covers both synthetic HPC pathologies rather than only the sharpest one. That test also shows the cross-algorithm comparison to be conservative in the direction that matters: pyABC's own penalty grows only to $4.8\times$, because it discards latecomers instead of waiting for the slowest draw (item iv), so the barrier's unmitigated cost exceeds the gap we report against it. The twin has since also been run at scale on the Cellular Potts workload (Table~\ref{tab:twin-cpm}), and there the answer is more qualified. The barrier costs $1.9\times$ from one to four nodes and $2.8\times$ at eight, while pyABC's penalty over the same range grows from $2.4\times$ to $5.1\times$: the barrier is the largest single component we can identify, but it is not the whole gap, and the residual --- pyABC's probabilistic-rejection step, its classical population weighting, its own coordination --- grows with worker count. On the realistic simulator we therefore attribute the gain to synchronization only in part, and say so rather than generalizing the straggler result. The attribution now covers all three benchmarks on which we report systems gains, with the caveat that its strength differs across them: near-total under a straggler or heterogeneity, partial on Cellular Potts. Separately, the twin is a diagnostic rather than a robust implementation --- it intermittently deadlocks at $\ge192$ ranks (five hangs while collecting Table~\ref{tab:twin-cpm}), so replicates there needed retrying, though all twenty entries are complete runs. Separately, on a methodological note for reproducibility: the post-run retroactive estimator builds its denominator from $m\le S{+}1=21$ history-reconstructed proposals (with a defensive prior floor), not the full evaluated history; for a near-instantaneous simulator whose history reaches many millions of particles the pass is evaluated in bounded-memory chunks, which changes memory, not the estimate.

\section{Conclusion}\label{sec:conclusion}
We presented a generation-free, single-arrival-driven ABC algorithm that combines smooth-kernel ABC with streaming AMIS reweighting, reports an estimator that is a pure function of the evaluated history and is therefore recoverable by replay, and integrates into Propulate's barrier-free island model. The estimator it reports is consistent, and satisfies a central limit theorem with an explicit asymptotic variance under rate conditions we measure rather than assume --- finding that the configuration reported here does not meet them, so the central limit theorem covers damped-adaptation variants; both are proved directly rather than inherited, under conditions relating its denominator to the density its sample was drawn from, which we state precisely, decompose, and partly measure on production runs --- but do not establish for the configuration the experiments use. The reported posterior is validated empirically by simulation-based calibration --- well calibrated in low dimensions and under multimodality, and conservative rather than over-confident in higher dimensions, where we trace the residual miscalibration to how much of the evaluated history is reported (the efficient top-$k$ proposal would under-cover; the full history over-covers) rather than to the weighting. Against a synchronous pyABC baseline matched on the ABC kernel and bandwidth schedule, and scored on each method's own reported posterior against a reference posterior rather than on a point mass at the truth, the verdict splits by dimension --- the asynchronous method is roughly twice as accurate on the one-dimensional analytic target and roughly twice as inaccurate on the four-dimensional one, where its weights intermittently degenerate and accuracy tracks the effective sample size at a log--log correlation of $-0.80$ --- while the method delivers substantial HPC benefits on heterogeneous and straggler-prone workloads, including a realistic Cellular Potts simulation.

That comparison also turned up the result we would most want a reader to carry away, because it bounds the first one. The reported posterior's effective sample size is set by the archive size and not by the simulation budget: $237$--$303$ effective particles out of $1.1$--$1.3\times10^6$ evaluated, flat across a thirteen-fold growth of the history within a single run, with $\mathrm{ESS}/k$ between $1.9$ and $3.6$ across a controlled sweep of both. Reporting at the tightest bandwidth reached, while that bandwidth is chosen by ESS-retention bisection against the archive, leaves an effective support of a few multiples of $k$ no matter how many particles are evaluated --- and it does so exactly when the target is identified enough for the bandwidth to tighten, which is to say exactly when the posterior is worth having. Removing the barrier therefore buys evaluations, reliably and by one to two-and-a-half orders of magnitude on the workloads we test; converting those evaluations into posterior accuracy is a separate problem, and on this evidence it is the one that now limits the method. We think generation-free ABC is the right execution model for large-scale simulator-based inference on modern HPC systems, and that the next work on it belongs in the reporting rule rather than in the scheduler.

